\begin{itemize}[leftmargin=*,label={}]
  \item \textbf{Progress\~ao aritm\'etica.}
  \[
    1+2+\cdots+n = \frac{n(n+1)}{2}
  \]
  \[
    a + (a+d) + \cdots + (a+(n-1)d)
    = \frac{n(2a+(n-1)d)}{2}
  \]

  \item \textbf{Progress\~ao geom\'etrica.}
  \[
    1+r+r^2+\cdots+r^n = \frac{r^{n+1}-1}{r-1}\qquad (r\ne1)
  \]
  \[
    a+ar+\cdots+ar^{n-1} = a\frac{r^n-1}{r-1}\qquad (r\ne1)
  \]

  \item \textbf{Somas de pot\^encias comuns.}
  \[
    \sum_{i=1}^n i = \frac{n(n+1)}{2}
  \]
  \[
    \sum_{i=1}^n i^2 = \frac{n(n+1)(2n+1)}{6}
  \]
  \[
    \sum_{i=1}^n i^3 = \left(\frac{n(n+1)}{2}\right)^2
  \]

  \item \textbf{Truques de floor/ceiling.}
  \[
    \left\lceil \frac{a}{b} \right\rceil = \left\lfloor \frac{a+b-1}{b} \right\rfloor
    \qquad (a,b>0)
  \]
  \[
    \left\lceil \frac{a}{b} \right\rceil = \frac{a+b-1}{b}
  \]
  em aritm\'etica inteira, para $a,b>0$.

  \ifextended
  \item \textbf{Coeficiente binomial.}
  \[
    \binom{n}{k}=\frac{n!}{k!(n-k)!}
  \]
  \fi

  \item \textbf{Coeficiente de multiconjunto.}
  \[
    \left(\!\binom{n}{k}\!\right)=\binom{n+k-1}{k}
  \]
  N\'umero de formas de escolher $k$ elementos entre $n$ tipos, com repeti\c{c}\~ao.

  \item \textbf{Estrelas e Barras.}
  O n\'umero de solu\c{c}\~oes inteiras n\~ao negativas de
  \[
    x_1+\cdots+x_k=n
  \]
  \'e
  \[
    \binom{n+k-1}{k-1}.
  \]
  O n\'umero de solu\c{c}\~oes inteiras positivas \'e
  \[
    \binom{n-1}{k-1}.
  \]

  \item \textbf{N\'umero de Catalan.}
  \[
    C_n=\frac{1}{n+1}\binom{2n}{n}
       =\binom{2n}{n}-\binom{2n}{n+1}
  \]

  \item \textbf{Desarranjos.}
  \[
    D_n = n!\sum_{k=0}^{n}\frac{(-1)^k}{k!}
  \]
  e
  \[
    D_n = (n-1)(D_{n-1}+D_{n-2})
  \]

  \item \textbf{Contagem / soma de divisores.}
  Se
  \[
    n=\prod p_i^{e_i},
  \]
  ent\~ao
  \[
    d(n)=\prod (e_i+1),
  \]
  \[
    \sigma(n)=\prod \frac{p_i^{e_i+1}-1}{p_i-1}.
  \]

  \ifextended
  \item \textbf{Phi de Euler.}
  \[
    \varphi(n)=n\prod_{p\mid n}\left(1-\frac1p\right)
  \]
  \fi

  \item \textbf{F\'ormula de Legendre.}
  \[
    v_p(n!)=\sum_{k\ge1}\left\lfloor \frac{n}{p^k}\right\rfloor
  \]

  \ifextended
  \item \textbf{Hash polinomial de substring.}
  Se
  \[
    H[i+1]=H[i]\cdot B+s[i],
  \]
  ent\~ao
  \[
    hash(l,r)=H[r]-H[l]\cdot B^{r-l}
  \]
  para o intervalo semiaberto $[l,r)$.
  \fi

  \item \textbf{Transforma\c{c}\~oes da dist\^ancia Manhattan.}
  Em 2D,
  \ifextended
  \[
    |x_1-x_2|+|y_1-y_2|
    = \max\Bigl(
      |(x+y)'-(x+y)|,\,
      |(x-y)'-(x-y)|
    \Bigr)
  \]
  Na pr\'atica, para a maior dist\^ancia entre pares:
  \[
    \max_{i,j}|x_i-x_j|+|y_i-y_j| = \max(D_+, D_-)
  \]
  onde
  \[
    D_+ = \max(x+y)-\min(x+y),\qquad
    D_- = \max(x-y)-\min(x-y)
  \]
  \else
  \[
    L_1 = \max(|\Delta_+|, |\Delta_-|)
  \]
  \[
    \Delta_\pm = (x\pm y)'-(x\pm y)
  \]
  Na pr\'atica:
  \[
    \max L_1 = \max(D_+, D_-)
  \]
  \[
    D_\pm = \max(x\pm y)-\min(x\pm y)
  \]
  \fi

  \ifextended
  \item \textbf{Dist\^ancia de ponto a reta.}
  Dist\^ancia do ponto $P$ \`a reta passando por $A,B$:
  \[
    \frac{|\operatorname{cross}(B-A,P-A)|}{|B-A|}
  \]
  \fi

  \item \textbf{F\'ormula de Heron.}
  Para lados $a,b,c$ e semiper\'imetro
  \[
    s=\frac{a+b+c}{2},
  \]
  a \'area do tri\^angulo \'e
  \[
    \sqrt{s(s-a)(s-b)(s-c)}
  \]

  \ifextended
  \item \textbf{Lei dos cossenos.}
  \[
    c^2=a^2+b^2-2ab\cos C
  \]
  \fi

  \ifextended
  \item \textbf{\^Angulo pelo produto escalar.}
  \[
    \cos\theta = \frac{a\cdot b}{|a||b|}
  \]
  \fi

  \item \textbf{Raio da circunfer\^encia circunscrita e inscrita.}
  Para \'area do tri\^angulo $\Delta$:
  \[
    R=\frac{abc}{4\Delta},\qquad r=\frac{\Delta}{s}
  \]
  onde $s=\frac{a+b+c}{2}$.

  \ifextended
  \item \textbf{Valor esperado por indicadores.}
  Se
  \[
    X=\sum_i I_i,
  \]
  ent\~ao
  \[
    \mathbb{E}[X]=\sum_i \Pr(I_i=1)
  \]
  \fi

  \ifextended
  \item \textbf{Contagem de itera\c{c}\~oes em submasks.}
  N\'umero total de pares $(m,s)$ com $s\subseteq m\subseteq \{0,\dots,n-1\}$:
  \[
    3^n
  \]
  \fi

  \item \textbf{Substrings distintas.}
  Usando suffix array:
  \[
    \#\text{substrings distintas}
    = \frac{n(n+1)}{2} - \sum lcp[i]
  \]
  Usando suffix automaton:
  \begin{center}
    \resizebox{0.95\linewidth}{!}{$\displaystyle
      \#\text{substrings distintas}
      = \sum_{v\neq root} \left(len(v)-len(link(v))\right)
    $}
  \end{center}

  \ifextended
  \item \textbf{Lembrete de interpreta\c{c}\~oes de Catalan.}
  \[
    C_n
  \]
  conta:
  sequ\^encias balanceadas de par\^enteses de tamanho $2n$, formas de \'arvores bin\'arias enraizadas com $n$ n\'os internos, triangula\c{c}\~oes de um $(n+2)$-gono, matchings perfeitos sem cruzamento em $2n$ pontos.
  \fi
\end{itemize}

\ifextended

\begin{itemize}[leftmargin=*,label={}]
  \item \textbf{N\'umeros de Stirling do segundo tipo.}
  \[
    S(n,k)=S(n-1,k-1)+k\,S(n-1,k)
  \]
  Forma fechada:
  \[
    S(n,k)=\frac1{k!}\sum_{i=0}^{k}(-1)^i\binom{k}{i}(k-i)^n
  \]

  \item \textbf{N\'umeros de Stirling do primeiro tipo (sem sinal).}
  \[
    c(n,k)=c(n-1,k-1)+(n-1)c(n-1,k)
  \]

  \item \textbf{N\'umeros de Bell.}
  \[
    B_n=\sum_{k=0}^{n}S(n,k)
  \]

  \item \textbf{Triangula\c{c}\~oes de pol\'igono convexo.}
  N\'umero de triangula\c{c}\~oes de um $n$-gono:
  \[
    C_{n-2}=\frac{1}{n-1}\binom{2n-4}{n-2}
  \]

  \item \textbf{Comprimento da corda.}
  Para raio $r$ e \^angulo central $\theta$:
  \[
    2r\sin\frac{\theta}{2}
  \]

  \item \textbf{Calota esf\'erica / f\'ormulas avan\c{c}adas de geometria.}
  Em geral n\~ao valem o espa\c{c}o no notebook, a menos que sua regi\~ao realmente cobre isso.
\end{itemize}

\fi
